\section{报告范围与结论}

\begin{frame}{如何阅读这份报告}
\begin{columns}[T,onlytextwidth]
\column{0.31\textwidth}
\begin{block}{第一层：事实}
当前 \pathcode{backend/} 源码、配置、测试、Mongo 和 trace；所有“已实现”结论可回到文件定位。
\end{block}
\column{0.31\textwidth}
\begin{block}{第二层：解释}
把控制流、状态、公式和数据口径连起来，解释为何系统会产生当前运行结果。
\end{block}
\column{0.31\textwidth}
\begin{block}{第三层：设计}
DMAIC 与论文实验是未来方案，使用 \proposal/\research 标识，不冒充已有能力。
\end{block}
\end{columns}

\vspace{0.5em}
\begin{center}
\impl\ 当前代码 \quad
\obs\ 真实运行 \quad
\risk\ 已识别限制 \quad
\proposal\ 工程改进 \quad
\research\ 待验证命题
\end{center}

\smallsource{完整证据索引：\pathcode{report/code_evidence_matrix.md}；运行快照：\pathcode{report/runtime_snapshot.json}}
\end{frame}

\begin{frame}{执行摘要：系统已可运行，但“有效覆盖”尚未闭环}
\begin{columns}[T,onlytextwidth]
\column{0.32\textwidth}
\begin{exampleblock}{已经成立}
\begin{itemize}
  \item 四个真实源与统一 DTO
  \item SDK tool loop + 本地 MCP 白名单
  \item 相关性三层管线
  \item Mongo 全局去重、trace 与 item-KG 交接
  \item adaptive 的 CorpusGap、UCB、checkpoint 与 playbook
\end{itemize}
\end{exampleblock}
\column{0.32\textwidth}
\begin{alertblock}{当前断点}
\begin{itemize}
  \item full/bootstrap 仍是 legacy 配额路径
  \item 最新 full 仅使用 NVD，15 个 Core 配额未满足
  \item 67.4\% item 为 uncertain，MSA 未完成
  \item arXiv 占全库 75.0\%，来源结构偏斜
  \item summary 字段可被误读为“零 gap”
\end{itemize}
\end{alertblock}
\column{0.32\textwidth}
\begin{block}{最优先方向}
\begin{enumerate}
  \item 先校准测量系统
  \item 把 full 做成可恢复的 adaptive campaign
  \item 用有效证据而非 item 数作为北极星
  \item 建立冻结 benchmark 与同预算 A/B
  \item 让全局 KG 结构空洞反向驱动采集
\end{enumerate}
\end{block}
\end{columns}
\end{frame}

\begin{frame}{证据基线：同一时间点、四种证据相互校验}
\centering
\metric{6,949}{Mongo 全局去重 items}\hfill
\metric{12}{持久化 runs}\hfill
\metric{86/86}{测试通过}\hfill
\metric{137}{最新 trace 事件}

\vspace{1.2em}
\begin{columns}[T,onlytextwidth]
\column{0.48\textwidth}
\begin{block}{静态证据}
\begin{itemize}
  \item 生产代码、schema、环境变量与容器编排
  \item 单元/集成测试覆盖的实际不变量
  \item 当前 topic taxonomy 与数据源参数
\end{itemize}
\end{block}
\column{0.48\textwidth}
\begin{block}{动态证据}
\begin{itemize}
  \item Mongo：跨 run 总量、来源、主题、标签与方法
  \item latest run：预算、动作、未满足配额
  \item JSONL trace：模型回合、工具、成本与降级
\end{itemize}
\end{block}
\end{columns}

\smallsource{快照时间 2026-07-15 12:48:10 UTC；测试同日执行。Mongo 是当前权威主存储。}
\end{frame}

\begin{frame}{显式需求覆盖矩阵}
\small
\begin{tabularx}{\textwidth}{>{\bfseries}p{0.19\textwidth} X p{0.25\textwidth}}
\toprule
用户要求 & 报告内容 & 主要证据 \\
\midrule
架构 & 包边界、部署、控制面/数据面、状态与数据流 & controller、tools、stores \\
原理 & ReAct 式闭环、混合控制、相关性门控、证据建模 & loop、hooks、relevance \\
算法/数理 & quota、有效证据、freshness、UCB、reward、emergence & gaps.py、bandit.py \\
代码对应 & 每个核心公式与行为的文件/行号映射 & code evidence matrix \\
运行结果 & 86 tests、smoke、Mongo 总量、latest run、trace 成本 & runtime snapshot \\
未来展望 & 从测量到可恢复 full campaign、全局图闭环 & DMAIC 9 页设计 \\
学术化 & 范围综述、研究空白、创新、RQ、实验与统计 & literature review + plan \\
\bottomrule
\end{tabularx}
\end{frame}

\begin{frame}{本报告的核心判断}
\begin{center}
\begin{minipage}{0.88\textwidth}
\Large
这个项目真正有研究价值的部分，不是“让 LLM 自己搜索”，而是把
\textcolor{Blue}{异质来源、证据质量、覆盖结构、预算约束、停止风险与可恢复性}
统一到一个可测、可控、可复现实验的主动采集系统中。
\end{minipage}
\end{center}

\vspace{1.2em}
\begin{block}{因此，评价函数必须从 volume 转向 value}
\scriptsize
\[
\begin{aligned}
U_{intel}&\neq \#\text{items},\\
U_{intel}&=f(\text{precision},\text{coverage},\text{independence},\text{freshness},\text{cost}).
\end{aligned}
\]
\end{block}
\end{frame}
